% Generated draft: requires mathematical and source review.
\documentclass[11pt]{article}
\usepackage[margin=25mm]{geometry}
\usepackage{fontspec}
\setmainfont{Latin Modern Roman}
\usepackage{amsmath,amssymb,mathtools,unicode-math}
\setmathfont{Latin Modern Math}
\usepackage{enumitem,needspace,xurl,hyperref}
\setlength{\parindent}{0pt}
\setlength{\parskip}{5pt}
\setlength{\emergencystretch}{3em}
\newcommand{\N}{\mathbb{N}}
\newcommand{\Z}{\mathbb{Z}}
\newcommand{\Q}{\mathbb{Q}}
\newcommand{\R}{\mathbb{R}}
\newcommand{\C}{\mathbb{C}}
\newcommand{\F}{\mathbb{F}}
\newcommand{\A}{\mathbb{A}}
\newcommand{\PP}{\mathbb{P}}
\newcommand{\E}{\mathbb{E}}
\newcommand{\eps}{\varepsilon}
\newcommand{\dd}{\,\mathrm{d}}
\begin{document}
\section*{UnsolvedMath ID: 30002106}
\textbf{Source problem code:} catalogue entry
\subsection{Definitions and mathematical statement}
The supplied record concerns an irreducible cuspidal representation \(\pi_v\) of \(\operatorname{GL}_g(F_v)\), subject to the hypothesis that its reduction modulo \(\ell\) is supercuspidal. Here \(\operatorname{GL}_g(F_v)\) denotes the group of invertible \(g\times g\) matrices over \(F_v\), and \(\overline{\mathbb F}_\ell\) denotes an algebraic closure of the field with \(\ell\) elements, for a prime \(\ell\). The field \(F_v\), the coefficient field of \(\pi_v\), and the hypotheses and lattice needed to define its reduction are not specified in the supplied record. The meanings of cuspidality, supercuspidality, and reduction must therefore be fixed in the intended representation-theoretic setting before the question is fully determined.

The integer parameters \(d\) and \(g\) occur in the range and shift below; interpreting the display requires positive integers \(d,g\), but their geometric roles and any further restrictions are not supplied. For each integer \(t\) with \(1\leq t\leq\lfloor d/g\rfloor\), the notation \(\mathcal F_{\overline{\mathbb F}_\ell}(\pi_v,t)\) denotes the object named in the source. Its construction, underlying space or stratum, and dependence on auxiliary choices are unspecified. The notation \([n]\) denotes cohomological shift by \(n\). The operators \({}^{p}j_{!*}\) and \({}^{p+}j_{!*}\) are presented as intermediate-extension functors for structures labelled \(p\) and \(p+\), respectively; neither the immersion \(j\), the ambient derived category, nor these two structures is defined in the record.

\textbf{Question as supplied.} For every irreducible cuspidal representation \(\pi_v\) of \(\operatorname{GL}_g(F_v)\) whose reduction modulo \(\ell\) is supercuspidal, does one have, for every integer \(t\) satisfying \(1\leq t\leq\lfloor d/g\rfloor\), an isomorphism
\[
{}^{p}j_{!*}\mathcal F_{\overline{\mathbb F}_\ell}(\pi_v,t)[d-tg]
\;\simeq\;
{}^{p+}j_{!*}\mathcal F_{\overline{\mathbb F}_\ell}(\pi_v,t)[d-tg] \, ?
\]
The shift placement and coefficient notation are retained from the supplied statement. If \(d<g\), the displayed range contains no integers, so this universal assertion is vacuous; the record does not specify whether that case is intended to be excluded.

\textbf{Unresolved qualifications.} The title places the question at parahoric level, but no parahoric subgroup, Shimura variety, or geometric hypotheses are specified. In addition, the use of two structures \(p\) and \(p+\) together with the displayed field coefficients needs clarification: the record does not say whether an integral coefficient category or a reduction of integral intermediate extensions is intended. These choices cannot be supplied by silently changing the formula. Thus the display is a faithful transcription of a conjectural question, conditional on recovering its missing setting, rather than a fully specified standalone conjecture.
\subsection{Short English statement}
The question asks whether the two indicated intermediate extensions agree for every allowed \(t\), assuming that \(\pi_v\) is irreducible cuspidal and its reduction modulo \(\ell\) is supercuspidal. The supplied record leaves the geometry, coefficient conventions, and the distinction between \(p\) and \(p+\) unspecified. Its label \texttt{partially\_solved} is supplied dated curation, not an independently verified conclusion about this formula.
\subsection{Sources}
\begin{itemize}
\item Local UnsolvedMath registry, permanent ID \(30002106\), source problem code \texttt{catalogue entry}, title \emph{Perverse Intermediate Extensions at Parahoric Level}; supplied complete source record associated with \texttt{data/unsolvedmath/problems.json} and the supplemental registry \texttt{catalog/problems.json}. The record lists proposed year 2012 and supplies no named proposer.
\item Supplied report citation: \emph{Reductions of Shimura Varieties (2013)}, Oberwolfach Reports, DOI 10.4171/owr/2012/32, pp.\ 1961--2011. \url{https://doi.org/10.4171/owr/2012/32}. The background describes Oberwolfach Reports as documenting mathematical workshops and their open-problem discussions. The title year 2013 and proposed year 2012 are retained as supplied.
\item Supplied literature source: Pappas and Zhu, \emph{local models and Kottwitz's conjecture}. \url{https://arxiv.org/abs/1110.5588}.
\item Supplied literature source: Gaitsgory and Haines, \emph{nearby cycles}. \url{https://arxiv.org/abs/math/0103047}. This attribution and descriptive label are reproduced from the record without independent verification.
\item Supplied literature source: \emph{original OWR report}. \url{https://doi.org/10.4171/owr/2012/32}.
\item Supplied literature assessment dated exactly \texttt{2026-08-21}: status \texttt{partially\_solved}. The assessment says that the relevant nearby-cycles and local-model conjectural framework has been proved in important cases, including character-type representations and several parahoric settings, while the extracted record does not support a general proof for arbitrary local representation. This is supplied dated curation derived from the OWR material, not independently verified here; the background explicitly requests independent verification of time-sensitive status and cited claims. It does not identify a precise theorem establishing a case of the displayed isomorphism.
\item The registry labels the statement \texttt{corrected\_verified} and reports that its mathematical content was checked against the official report, with notation, context, and question form normalized without changing the source-level claim. This verification claim is supplied metadata, not a source consultation performed for this draft.
\end{itemize}
\end{document}
